% ============================================================================
% Figure 3: Combined Meta-Learning Cost Characterisation (2-panel)
% ============================================================================
% Unified story: meta-learning overhead when priors are strong,
% with weight evolution explaining the mechanism.
% Updated 2026-02-14 with corrected results (N=20, η=0.030).

\begin{figure}[t]
    \centering
    \includegraphics[width=\textwidth]{experiments_v1/03_figure/results/figure3_combined.pdf}
    \caption{%
        \textbf{The Cost of Adaptive Safety.}
        %
        \textbf{(A) Strategy Comparison:}
        When warmup priors are strong, direct prior usage (warmup-only)
        achieves 30.0~$\pm$~2.0 cumulative regret---substantially lower
        than Corralling meta-learning (48.2~$\pm$~5.1, +60\% overhead).
        Tabula rasa (49.2~$\pm$~3.4) provides the prior-free baseline.
        That Corralling~$\approx$~tabula rasa confirms the meta-learning
        overhead roughly equals the prior advantage on this benchmark.
        Error bars: $\pm$1 std across $N{=}20$ seeds.
        %
        \textbf{(B) Weight Evolution Explains the Gap:}
        Corralling correctly learns to allocate 70.6\%~$\pm$~11.0\%
        weight to the warmup expert over 750 queries, validating that
        meta-learning identifies the superior strategy.
        However, the gradual convergence from 50\% to 71\% warmup weight
        accumulates the 18-regret overhead observed in~(A), as the
        meta-learner cannot fully exploit its growing confidence within
        the Exp4 worst-case framework.
        Individual seed trajectories (thin lines) and mean trajectory
        (bold) with confidence bands show smooth, reproducible
        convergence with low cross-seed variance.
        %
        \textbf{Key insight:}
        Meta-learning is not free.  The 60\% overhead is the insurance
        premium paid for adaptive safety in a benign environment.
        Sections~\ref{sec:catastrophic} and~\ref{sec:zero_shot}
        demonstrate when this premium pays off: under model failure and
        zero-shot adoption, Corralling detects and adapts within
        3--50 steps---capability that warmup-only lacks entirely.
        %
        Setup: 750-prompt LMSYS holdout, $N{=}20$ seeds, sampling
        without replacement.
        All configurations use $\eta{=}\sqrt{\ln 2/750}{\approx}0.030$,
        $\gamma{=}0.05$, $\alpha{=}2.0$
        (see Appendix~\ref{app:ablations} for ablation studies).
    }
    \label{fig:metalearning_cost}
\end{figure}

% ============================================================================
% Alternative: Separate figures if page budget allows
% ============================================================================
\iffalse
% Figure 3A: Convergence Comparison Only
\begin{figure}[t]
    \centering
    \includegraphics[width=0.65\columnwidth]{experiments_v1/03_figure/results/convergence/figure_convergence_dynamics.pdf}
    \caption{%
        \textbf{The Cost of Adaptive Safety.}
        Corralling incurs 60\% higher cumulative regret
        (48.2~$\pm$~5.1) compared to directly using warmup priors
        (30.0~$\pm$~2.0).
        Tabula rasa (49.2~$\pm$~3.4) performs comparably to Corralling,
        confirming the meta-learning overhead roughly equals the prior
        advantage.
        This 60\% premium is the cost of hedging against prior failure
        in a benign environment---an insurance policy that pays off
        under model degradation (\S\ref{sec:catastrophic}).
    }
    \label{fig:convergence_comparison}
\end{figure}

% Figure 3B: Weight Evolution Only
\begin{figure}[t]
    \centering
    \includegraphics[width=0.85\columnwidth]{experiments_v1/03_figure/results/weight_evolution/figure_weight_evolution.pdf}
    \caption{%
        \textbf{Weight Evolution Explains Meta-Learning Overhead.}
        \textbf{(A)} Individual seed trajectories show Corralling
        learning to trust the warmup expert over 750 queries.
        \textbf{(B)} Mean trajectory converges to 70.6\% warmup weight
        (grey dashed line shows initial 50\% equal weighting).
        The learning period from 50\% to 71\% allocation accumulates the
        18 regret overhead.
        The 11\% standard deviation in final weights indicates stable
        convergence across data orderings, confirming the theoretically
        optimal $\eta$ produces reproducible behaviour.
    }
    \label{fig:weight_evolution}
\end{figure}
\fi
